\chapter{Linear Systems}

\section{Triangular Systems}

\begin{exercise}[Forward Substitution]
    Write a function that performs forward substitution on a $n\times n$ lower triangular matrix.
\end{exercise}

\begin{exercise}
    Test your code on the linear systems:
    $$
        \mathbf{a)}
        \quad 
        \begin{bmatrix}
            -2 & 0 & 0 \\
            1 & -1 & 0 \\
            3 & 2 & 1 
        \end{bmatrix} 
        \mathbf{x} = 
        \begin{bmatrix}
            -4 \\ 2 \\ 1
        \end{bmatrix};
        \qquad
        \mathbf{b)} 
        \quad 
        \begin{bmatrix}
            4 & 0 & 0 & 0 \\
            1 & -2 & 0 & 0 \\
            -1 & 4 & 4 & 0 \\
            2 & -5 & 5 & 1
        \end{bmatrix} 
        \mathbf{x} = 
        \begin{bmatrix}
            -4 \\ 1 \\ -3 \\ 5
        \end{bmatrix}
    $$
    You can verify the solution by solving the system by hand and/or computing $\mathbf{L}\mathbf{x}$ and subtracting $\mathbf{b}$.
\end{exercise}

\begin{exercise}
    Write a function that performs backward substitution on a $n\times n$ upper triangular matrix.
\end{exercise}

\begin{exercise}
    Test your code on the linear systems:
    $$
    \begin{bmatrix}
      3 & 1 & 0  \\
      0 & -1 & -2  \\
      0 & 0 & 3  \\
    \end{bmatrix} \mathbf{x} = \begin{bmatrix}
      1 \\ 1 \\ 6
    \end{bmatrix}
    \qquad
    \begin{bmatrix}
      3 & 1 & 0 & 6 \\
      0 & -1 & -2 & 7 \\
      0 & 0 & 3 & 4 \\
      0 & 0 & 0 & 5
    \end{bmatrix} \mathbf{x} = \begin{bmatrix}
      4 \\ 1 \\ 1 \\ 5
    \end{bmatrix}
    $$
    You can verify the solution by solving the system by hand and/or computing $\mathbf{U}\mathbf{x}$ and subtracting $\mathbf{b}$.
\end{exercise}

\section{LU Decomposition}

\begin{exercise}
    Write a function that performs the LU-decomposition of a $n\times n$ matrix. 
\end{exercise}

\begin{exercise}\label{LU Factorization}
    For each matrix, compute their LU-factorization and verify the correctness.

    $$
        (a) \, \begin{bmatrix}
            2 & 3 & 4 \\
            4 & 5 & 10 \\
            4 & 8 & 2
        \end{bmatrix}\qquad
        (b) \, \begin{bmatrix}
            6 & -2 & -4 & 4\\
            3 & -3 & -6 & 1 \\
            -12 & 8 & 21 & -8 \\
            -6 & 0 & -10 & 7
        \end{bmatrix}\qquad
        (c) \, \begin{bmatrix}
            1 & 4 & 5 & -5 \\
            -1 & 0 & -1 & -5 \\
            1 & 3 & -1 & 2 \\
            1 & -1 & 5 & -1 
        \end{bmatrix}
    $$
\end{exercise}

\begin{exercise}
    The matrices
    $$
    \mathbf{T}(x,y) = \begin{bmatrix}
        1 & 0 & x \\ 
        0 & 1 & y \\ 
        0 & 0 & 1
    \end{bmatrix},\qquad
    \mathbf{R}(\theta) = \begin{bmatrix}
        \cos\theta & \sin \theta & 0 \\ 
        -\sin\theta & \cos \theta & 0 
        \\ 0 & 0 & 1
    \end{bmatrix}
    $$
    are used to represent translations and rotations of plane points in computer graphics. For the following, let
    $$
        \mathbf{A} = \mathbf{T}(3,-1)\mathbf{R}(\pi/5)\mathbf{T}(-3,1), 
        \qquad
        \mathbf{z} = 
        \begin{bmatrix}
            2 \\ 2 \\ 1
        \end{bmatrix}.
    $$
    \begin{enumerate}
        \item Find $\mathbf{b} = \mathbf{A}\mathbf{z}$.
        \item Find the LU factorization of $\mathbf{A}$.
        \item Use the factors with triangular substitutions in order to solve $\mathbf{A}\mathbf{x}=\mathbf{b}$, and find $\mathbf{x}-\mathbf{z}$.
    \end{enumerate}
\end{exercise}

\begin{exercise}
    Compute the determinant of the matrices in Exercise \ref{LU Factorization} using their LU-decomposition.
\end{exercise}

\begin{exercise}
    Write a function that performs the LU-decomposition of a $n\times n$ matrix including row pivoting. Apply it to the matrices (a)-(c) of Exercise \ref{LU Factorization}.
\end{exercise}

\section{Conditioning of Linear Systems}

\begin{exercise} 
    Let us consider the matrix:

    $$
        \mathbf{A} =
        \begin{bmatrix}
        -\epsilon & 1 \\ 
        1 & -1
        \end{bmatrix}.
    $$

    If $\epsilon=0$, LU factorization without partial pivoting fails for $\mathbf{A}$. But if $\epsilon\neq 0$, we can proceed without pivoting, at least in principle.

    \begin{enumerate}
        \item Construct $\mathbf{b}=\mathbf{Ax}$ taking $\epsilon=10^{-12}$ for the matrix $\mathbf{A}$ and $\mathbf{x}=[1,1]$ for solution.
        \item Factor the matrix using the LU-decomposition without pivoting and solve numerically for $\mathbf{x}$. How accurate is the result?
        \item Try to verify the LU-factorization by computing $\mathbf{A}-\mathbf{LU}$. Does it work?
        \item Repeat for $\epsilon=10^{-20}$. How accurate is the result now?
        \item Compute the condition number of the matrix $\mathbf{A}$; you can choose your favorite norm. Is the system $\mathbf{Ax}=\mathbf{b}$ ill-conditioned?
        % \item Compute the matrices $\mathbf{L}$ and $\mathbf{U}$ of the LU-factorization of $\mathbf{A}=\mathbf{LU}$ with *pen and paper*.
        \item Compute the condition numbers of $\mathbf{L}$ and $\mathbf{U}$. Are the triangular systems $\mathbf{L}\mathbf{z}=\mathbf{b}$ and $\mathbf{U}\mathbf{x}=\mathbf{z}$ ill-conditioned?
        \item Repeat the steps above for $\epsilon=10^{-20}$ for the matrix
        $$
        \mathbf{PA}=
            \begin{bmatrix}
            1 & -1 \\
            -\epsilon & 1 
            \end{bmatrix}.
        $$
        where we permuted the first and second row of $\mathbf{A}$. Note that this is the matrix that would be factorized when using row pivoting, since $\epsilon\ll1$.
    \end{enumerate}

\end{exercise}

\section{Overdetermined Linear Systems}


